\documentclass[a4paper,12pt]{article}

\usepackage[T2A]{fontenc}
\usepackage[utf8]{inputenc}
\usepackage[russian]{babel}
\usepackage{amsmath,amssymb,amsthm}
\usepackage{mathtools}
\usepackage{booktabs,array}
\usepackage{forest}
\usepackage[hidelinks]{hyperref}
\usepackage{tikz}
\usetikzlibrary{arrows.meta,positioning}

\title{Лабораторная работа №1}
\author{Лимонов Дмитрий ИУ9-52Б}

\begin{document}
\maketitle

\section*{Изначальная система переписывания}

Рассматривается система переписывания строк (вариант 15), где алфавит
\[
\Sigma = \{a,b,c\},
\]
а множество правил задано следующим образом:
\[
\begin{aligned}
\texttt{abca} &\;\to\; \texttt{aaaa},\\
\texttt{bcab} &\;\to\; \texttt{bbbb},\\
\texttt{cabc} &\;\to\; \texttt{cccc},\\
\texttt{abc}  &\;\to\; \texttt{aaa},\\
\texttt{abc}  &\;\to\; \texttt{bbb},\\
\texttt{abc}  &\;\to\; \texttt{ccc},\\
\texttt{aabbcc} &\;\to\; \texttt{abcabc}.
\end{aligned}
\]

\section*{Завершимость}

Рассмотрим правила переписывания. Все правила, кроме последнего, могут быть применены только в том случае, если в строке есть подстрока из 3 различных символов, идущих подряд (\texttt{abc, bac, cab, ...}), при этом все правила как раз убирают эти различные символы. Последнее же правило может быть применено в случае, если встречается подстрока \texttt{bb}. Так как среди предыдущих правил нет ни одного способа пополнить количество именно \texttt{bb}, то, в какой-то момент, переписывания остановятся

\section*{Конечность}

Исходя из структуры правил можно заметить, что ни одно из них не уменьшает длины изначальной строки. Значит, как минимум, существуют классы эквивалентности по длине строки.


\section*{Конфлюэнтность}

Изначальная система не является конфлюэнтной ни локально, ни глобально. Это можно увидеть на примере строки 
\texttt{abc}:

\begin{center}
\begin{forest}
for tree={
  draw,
  rounded corners,
  align=center,
  s sep=10mm,
  l sep=12mm,
  edge={-stealth, thick}, 
}
[abc
  [aaa]
  [bbb]
  [ccc]
]
\end{forest}
\end{center}

Дополним систему по алгоритму Кнута-Бендикса. Сначала введём лексикографический порядок в систему и переупорядочим правила:

\[
\begin{aligned}
\texttt{abca} &\;\to\; \texttt{aaaa},\\
\texttt{bcab} &\;\to\; \texttt{bbbb},\\
\texttt{cccc} &\;\to\; \texttt{cabc},\\
\texttt{abc}  &\;\to\; \texttt{aaa},\\
\texttt{bbb}  &\;\to\; \texttt{abc},\\
\texttt{ccc}  &\;\to\; \texttt{abc},\\
\texttt{abcabc} &\;\to\; \texttt{aabbcc}.
\end{aligned}
\]

Далее, мною была написана программа по подбору новых правил. Исходный код программы можно найти \href{https://github.com/dm800/TFLcheck}{\textcolor{blue}{здесь}}. В результате выполнения программы стало ясно, что в системе при пополнении возникает бесконечное количество правил, поэтому возьмём приближённую эквивалентной систему (она будет полностью эквивалентна на строках длиной 7 и меньше.):

\[
\begin{alignedat}{3}
\texttt{bcab}    &\;\to\; \texttt{aaab}   & \qquad \texttt{caaa}    &\;\to\; \texttt{aaac}   & \qquad \texttt{abc}     &\;\to\; \texttt{aaa}    \\[2pt]
\texttt{bbb}     &\;\to\; \texttt{aaa}    & \qquad \texttt{ccc}     &\;\to\; \texttt{aaa}    & \qquad \texttt{aabbcc}  &\;\to\; \texttt{aaaaaa} \\[2pt]
\texttt{baaa}    &\;\to\; \texttt{aaab}   & \qquad \texttt{aaacc}   &\;\to\; \texttt{aaaab}  & \qquad \texttt{aaabab}  &\;\to\; \texttt{aaaaac} \\[2pt]
\texttt{aaabbc}  &\;\to\; \texttt{aaabaa} & \qquad \texttt{aaacab}  &\;\to\; \texttt{aaaaaa} & \qquad \texttt{aaacbc}  &\;\to\; \texttt{aaacaa} \\[2pt]
\texttt{aaaaabb} &\;\to\; \texttt{aaaaaac}& \qquad \texttt{aaabaac} &\;\to\; \texttt{aaaaaaa}& \qquad \texttt{aaabacc} &\;\to\; \texttt{aaabaab}\\[2pt]
\texttt{aaabbaa} &\;\to\; \texttt{aaaaaac}& \qquad \texttt{aaabbab} &\;\to\; \texttt{aaaaaaa}& \qquad \texttt{aaacacc} &\;\to\; \texttt{aaacaab}\\[2pt]
\texttt{aaacbab} &\;\to\; \texttt{aaacaac}& \qquad \texttt{aaacbbc} &\;\to\; \texttt{aaacbaa}& \qquad {}                &                         \\
\end{alignedat}
\]


\section*{Тестирование}
\subsection*{Фаззинг}
Исходный код программы для фаззинг-тестирования может быть найден в репозитории. Основная идея - сгенерировать случайное слово, далее привести его к нормальной форме в $\tau$. После этого уже изначальное слово и полученную на предыдущем шаге нормальную форму перевести в общее слово, но уже в $\tau'$. В моей программе, для простоты реализации, этим словом также будет нормальная форма, то есть:


\[
\begin{aligned}
\begin{tikzpicture}[>=Stealth]
  \node (w)   {$w$};
  \node (wp)  [below=1.6cm of w] {$w'$};                    
  \node (wpp) [below right=1.2cm and 2.4cm of wp] {$w''$}; 
  
  \draw[->] (w)  -- node[left] {$\tau$}  (wp);
  \draw[->] (wp) -- node[above] {$\tau'$} (wpp);
  \draw[->] (w)  to[bend left=15] node[above right] {$\tau'$} (wpp);
\end{tikzpicture}
\end{aligned}
\]


\subsection*{Метаморфное тестирование}
Исходный код программы для метаморфного тестирования также может быть найден в репозитории. В качестве инвариантов были выбраны две характеристики:

1. Длина слова 

2. $(\Delta A - \Delta C) mod 3= 0$ , где $\Delta A$ и $\Delta C$ - разность числа букв \texttt{a} или \texttt{c} соответственно.


\end{document}